%\section{Early Experiments on Countdown with Qwen2.5-7B}
\section{Self-Revision Design Choices}
\label{app:countdown}

 %\sk{These experiments used an early/different version of the self-revision pipelinewith different prompts, hyperparameters, and base model than SD-Zero.Will rerun + update results + writeup below  with the full SD-Zero pipeline on Qwen3-4B-Instruct so that the prompts and setup are consistent across all experiments in the paper.}

%The experiments below use a different base model and pipeline configuration than the main \ours{} experiments; we include them because the design principles they reveal corroborate those of \ours{}.
%This section reports early experiments exploring self-revision and
%distillation on the Countdown task using Qwen2.5-7B.  
%This section reports experiments on the Countdown task with Qwen2.5-7B that explore design questions related to self-revision training. The setup differs from the main SD-ZERO experiments in base model, pipeline configuration, and hyperparameters; we include them as supplementary evidence because several findings are consistent with the design principles underlying SD-ZERO. 
%Countdown is a constrained arithmetic reasoning problem in which the model must combine
%a given set of numbers using basic operations to reach a target value.
%The experiments address three questions: (1)~how does the choice of
%training data source affect in-distribution performance? (2)~does
%self-revision transfer out-of-distribution, and does it complement RL?
%(3)~what design principles emerge for a self-revision pipeline?  Several
%of the answers---on-policy supervision outperforms off-policy,
%correctness filtering is essential, self-revision complements
%RL---consistent with key design choices in SD-Zero.

We present complementary experiments on Countdown, a constrained arithmetic reasoning task, using Qwen2.5-7B. Although the setup differs from the main experiments (different base model, task domain, and pipeline configuration), three findings align with and corroborate the design principles of \ours{}: (1) on-policy self-revision data outperforms off-policy teacher data, (2) correctness filtering of revision traces is essential, and (3) self-revision training provides a stronger foundation for subsequent training than standard SFT does.

 
\subsection{Effect of Training Data Source}
\label{app:countdown-sources}
 
We compare several training data sources for supervised fine-tuning on
Countdown.  For each of the 8K original SFT questions, we sample $N$ solutions per question and
optionally filter for correctness. Table~\ref{tab:countdown-sources} reports the results.
 
\begin{itemize}
\item \textbf{LLaMA-70B}: $N\!=\!16$, pass@$16 = 81\%$, filtered
      $\rightarrow$ 6.5K pairs.
\item \textbf{Qwen2.5-7B} (self-generated): $N\!=\!16$,
      pass@$16 = 78.7\%$, filtered $\rightarrow$ 6.2K pairs.
\item \textbf{GPT-4o}: $N\!=\!5$, pass@$5 = 67.5\%$, filtered
      $\rightarrow$ 5.4K pairs.
\item \textbf{SFT on GRPO-Qwen2.5-7B data} (on-policy): GRPO reaches
      pass@$1 = 91\%$ after RL; we distill its generations into
      6.5--8K SFT pairs without filtering.
\item \textbf{Self-Revision}: Qwen2.5-7B critiques and refines the previous
      response attempt; unfiltered (6.5K) and filtered
      (5.9K) variants.
\end{itemize}
 

 
\begin{table}[h]
\centering
\vspace{4pt}
\small
\begin{tabular}{lccc}
\toprule
\textbf{Training Source} & \textbf{pass@1} & \textbf{pass@2} & \textbf{pass@4} \\
\midrule
No training & 0.410 & 0.538 & 0.639 \\
\midrule
LLaMA-70B (6.5K, filtered) & 0.605 & 0.725 & 0.814 \\
Qwen2.5-7B (6.2K, filtered) & 0.552 & 0.676 & 0.768 \\
GPT-4o (5.4K, filtered) & 0.617 & 0.711 & 0.775 \\
\midrule
SFT on GRPO data (8K, unfiltered) & 0.878 & 0.904 & 0.922 \\
SFT on GRPO data (6.5K, unfiltered) & 0.883 & 0.910 & 0.919 \\
\midrule
Self-Revision (6.5K, unfiltered) & 0.529 & 0.639 & 0.717 \\
Self-Revision (5.9K, filtered) & 0.630 & 0.737 & 0.810 \\
\bottomrule
\end{tabular}
\caption{%
  Countdown pass@$k$ after fine-tuning Qwen2.5-7B on different data
  sources.  On-policy sources (SFT on GRPO data, self-revision)
  consistently outperform off-policy ones (LLaMA-70B, GPT-4o).  Among
  methods that do not require a prior RL stage, filtered self-revision
  achieves the highest pass@1.}
\label{tab:countdown-sources}
\end{table}
 
\paragraph{Findings.}
Three patterns emerge.  First, \textbf{on-policy data dominates
off-policy data}: SFT on data generated by the GRPO-trained checkpoint
(${\sim}0.88$) substantially outperforms SFT on LLaMA-70B ($0.605$) or
GPT-4o ($0.617$) data, despite comparable data sizes and correctness
filtering.  Note that GRPO itself reaches ${\sim}0.91$ via RL alone;
the ${\sim}0.88$ figure reflects SFT distillation of its outputs.
Distribution match matters more than teacher quality---the on-policy
6.5K subset performs comparably to the 8K variant.  \ours{} also
uses the student's own reviser as the teacher.
 
Second, \textbf{self-revision outperforms first-attempt sampling}:
filtered self-revision (pass@1=$0.630$) exceeds the model's own
correctness-filtered first attempts (pass@1=$0.552$).  Both are on-policy, so
the gain isolates the value of the revision step itself.
 
Third, \textbf{correctness filtering is essential}: unfiltered
self-revision ($0.529$) underperforms the base model's first attempts,
while filtering raises it to $0.630$---a $19\%$ relative gain.  Noisy
revision traces actively hurt.  \ours{}'s Phase~1 also retains only
successful revisions (Section~2.1).
 
%\paragraph{Comparison with GRPO.}
%GRPO reaches ${\sim}0.91$ in-distribution while being more
%sample-efficient: it trains directly with a reward signal, whereas
%self-revision requires generating candidates, revising, and filtering to
%produce 5.9K pairs.  Even SFT on GRPO-generated data reaches
%${\sim}0.88$ without any RL at inference time.  GRPO also largely
%preserves OOD performance (Table~\ref{tab:countdown-ood}).
%Self-revision SFT, while the strongest non-RL method that does not
%require a prior RL stage, does not match RL in-distribution on
%Countdown.  This pointed to a need to rethink the pipeline.
% 
%However, two findings suggested self-revision remained a valuable
%\emph{component}.  First, it provides a better initialization for RL
%than off-policy data (Section~\ref{app:countdown-ood}).  Second, the
%design principles above---on-policy supervision, correctness filtering,
%revision over raw sampling---identified a clear path for improving the
%pipeline rather than abandoning it.  SD-Zero's two-phase procedure
%(Section~2) shares each of these properties.
% 
%The in-distribution lift from self-revision SFT here is comparable to
%that of SRT in the main experiments (Table~1), but the GRPO ceiling on
%Countdown is much higher ($0.91$ vs.\ the GRPO results in Table~1),
%making the remaining gap appear larger.  Several confounds make this gap
%hard to interpret: the initial model differs (Qwen2.5-7B vs.\
%Qwen3-4B-Instruct); the self-revision prompts and hyperparameters differ
%from SD-Zero; and the GRPO implementation follows a r1 style DAPO-style recipe (with partial reward for wrong answer but right format).
 
%\subsection{Out-of-Distribution Evaluation}
\subsection{Self-Revision as Initialization for RL}

\label{app:countdown-ood}
 
We evaluate whether self-revised targets (A$'$) improve OOD
generalization, especially when combined with RL.  We fine-tune
Qwen2.5-7B on 8K Countdown questions with either LLaMA-70B--generated
answers~(A) or self-revised answers~(A$'$), optionally followed by GRPO.
Neither A nor A$'$ are filtered for correctness.  Note that A uses
off-policy teacher answers; Table~\ref{tab:countdown-sources} shows this
is a conservative baseline, since off-policy data underperforms
on-policy data even in-distribution.  All evaluations use a 4096-token
generation budget.
 
\begin{table}[h]
\centering
\vspace{4pt}
\footnotesize
\setlength{\tabcolsep}{1.5pt}
\begin{tabular}{l cc cc cc cc cc cc cc}
\toprule
& \multicolumn{2}{c}{\textbf{AIME24}}
& \multicolumn{2}{c}{\textbf{AIME25}}
& \multicolumn{2}{c}{\textbf{Countdown}}
& \multicolumn{2}{c}{\textbf{IFBench}}
& \multicolumn{2}{c}{\textbf{IFEval}}
& \multicolumn{2}{c}{\textbf{MMLU}}
& \multicolumn{2}{c}{\textbf{MMLU-Pro}} \\
\midrule
\textbf{Method}
& @1 & @128 & @1 & @128 & @1 & @128
& @1 & @128 & @1 & @128 & @1 & @128 & @1 & @128 \\
\midrule
Base
  & .07 & .27 & .07 & .43 & .41 & .90
  & .27 & .58 & .71 & .90 & .74 & .91
  & .56 & .82 \\
GRPO (Countdown)
  & .08 & .30 & .08 & .40 & .91 & .98
  & .27 & .61 & .72 & .91 & .74 & .92
  & .56 & .84 \\
\midrule
Countdown (A)
  & .05 & .27 & .05 & .43 & .39 & .91
  & .23 & .54 & .52 & .85 & .69 & .96
  & .46 & .91 \\
Countdown (A) + GRPO
  & .06 & .23 & .07 & .40 & .92 & .98
  & .24 & .55 & .53 & .84 & .71 & .97
  & .50 & .93 \\
\midrule
Countdown (A$'$)
  & .07 & .30 & .06 & .47 & .40 & .88
  & .19 & .62 & .60 & .91 & .71 & .97
  & .50 & .94 \\
Countdown (A$'$) + GRPO
  & .07 & \textbf{.40} & .07 & \textbf{.47} & .91 & .98
  & .20 & .58 & .62 & .91 & .72 & .97
  & .52 & .93 \\
\bottomrule
\end{tabular}
\caption{%
  OOD evaluation (pass@1 / pass@128) after training on Countdown with
  LLaMA-70B answers~(A) or self-revised answers~(A$'$), optionally
  followed by GRPO.  A$'$+GRPO achieves the highest OOD scores.}
\label{tab:countdown-ood}
\end{table}

%\paragraph{Findings.}
The key result is that \textbf{self-revised targets provide a better
initialization for RL on OOD benchmarks.}  Countdown~(A$'$)+GRPO yields
pass@128 of $0.40$ on AIME24 and $0.47$ on AIME25, compared to $0.23$
and $0.40$ for Countdown~(A)+GRPO, and $0.30$ and $0.40$ for standalone
GRPO.  In-distribution Countdown performance is identical across all
GRPO variants ($0.91$/$0.98$).

%\paragraph{Comparison with GRPO.}
%GRPO reaches pass@1=0.91 in-distribution, outperforming filtered self-revision (pass@1=0.63). Self-revision SFT, while the strongest non-RL method that does not require a prior RL stage, does not close this gap on Countdown. The two setups differ in base model and training recipe, so direct comparison is limited. However, self-revision provides a stronger initialization for subsequent RL than off-policy data, suggesting the two approaches are complementary rather than competing.
%In-distribution, GRPO substantially outperforms even the best self-revision SFT variant (pass@1=0.91 vs. 0.63; Tables \ref{tab:countdown-ood} and \ref{tab:countdown-sources} respectively). The two setups differ in base model and training recipe, so direct comparison is limited. However, the OOD results above suggest the two approaches are complementary rather than competing.
 
 
%This comparison is conservative in two respects.  First, A uses
%LLaMA-70B answers (off-policy), not the model's own first attempts;
%Table~\ref{tab:countdown-sources} shows that self-revision outperforms
%even on-policy first attempts ($0.630$ vs.\ $0.552$), so the advantage
%of A$'$ likely reflects the revision process itself.  Second, neither A
%nor A$'$ were filtered for correctness; given the $19\%$ relative gain
%from filtering in Table~\ref{tab:countdown-sources}, the gap would
%presumably widen further.

Notably, this comparison is conservative: neither A nor A' were filtered for correctness, and A uses off-policy 
%LLaMA-70B
answers rather than the model's own attempts. Despite this, self-revised targets still yield a meaningfully stronger RL initialization on OOD benchmarks (Table~\ref{tab:countdown-ood}). This suggests that the revision process itself — even without correctness filtering — carries useful signal, and that the gains would likely be larger with filtering (Table~\ref{tab:countdown-sources} shows a 19\% relative gain from filtering in-distribution).
Moreover, these experiments use GRPO rather than \ours{}'s distillation objective as the second stage, suggesting the value of self-revision as initialization is not specific to a particular second-stage method.
 
%The second-stage algorithm also differs: these experiments use GRPO
%rather than the SDFT distillation objective in SD-Zero's Phase~2.  The
%complementarity of self-revision with RL thus holds under a different
%second-stage method.

More broadly, the fact that unfiltered self-revision provides useful signal raises the question of whether self-revision training can extend to settings without verifiable rewards, where correctness filtering is unavailable. We leave this to future work.
 
%\paragraph{Tulu3 data.}
%We ran parallel experiments with Tulu3-generated training data, but the
%Tulu3 solutions contain extended thinking tokens and no correctness
%filtering was applied.  The A$'$ vs.\ A improvements were noisier,
%and consistent with the findings in Appendix~\ref{app:opsd-thinking} on the
%difficulty of supervising thinking-mode traces and with the importance of
%filtering documented above.